\begin{theorem}
\label{theorem-bondal-van-den-Bergh}
Let $X$ be a quasi-compact and quasi-separated scheme. The category
$D_\QCoh(\mathcal{O}_X)$ can be generated by a single
perfect object. More precisely, there exists a perfect object
$P$ of $D(\mathcal{O}_X)$ such that for 
$E \in D_\QCoh(\mathcal{O}_X)$ the following are equivalent
\begin{enumerate}
\item $E = 0$, and
\item $\Hom_{D(\mathcal{O}_X)}(P[n], E) = 0$ for all $n \in \mathbf{Z}$.
\end{enumerate}
\end{theorem}

\begin{proof}
We will prove this using the induction principle of
Cohomology of Schemes, Lemma \ref{coherent-lemma-induction-principle}.

\medskip\noindent
If $X$ is affine, then $\mathcal{O}_X$ is a perfect generator.
This follows from Lemma \ref{lemma-affine-compare-bounded}.

\medskip\noindent
Assume that $X = U \cup V$ is an open covering with $U$ quasi-compact
such that the theorem holds for $U$ and $V$ is an affine open.
Let $P$ be a perfect object of $D(\mathcal{O}_U)$ which is a generator
for $D_\QCoh(\mathcal{O}_U)$. Using
Lemma \ref{lemma-direct-summand-of-a-restriction} we may
choose a perfect object
$Q$ of $D(\mathcal{O}_X)$ whose restriction to $U$ is a direct sum one
of whose summands is $P$. Say $V = \Spec(A)$. Let $Z = X \setminus U$.
This is a closed subset of $V$ with $V \setminus Z$ quasi-compact.
Choose $f_1, \ldots, f_r \in A$ such that
$Z = V(f_1, \ldots, f_r)$. Let $K \in D(\mathcal{O}_V)$ be the perfect
object corresponding to the Koszul complex on $f_1, \ldots, f_r$ over $A$.
Note that since $K$ is supported on $Z \subset V$ closed, the pushforward
$K' = R(V \to X)_*K$ is a perfect object of $D(\mathcal{O}_X)$ whose
restriction to $V$ is $K$ (see
Cohomology, Lemma \ref{cohomology-lemma-pushforward-perfect}).
We claim that $Q \oplus K'$ is a generator for
$D_\QCoh(\mathcal{O}_X)$.

\medskip\noindent
Let $E$ be an object of $D_\QCoh(\mathcal{O}_X)$ such that
there are no nontrivial maps from any shift of $Q \oplus K'$ into $E$.
By Cohomology, Lemma \ref{cohomology-lemma-pushforward-restriction}
we have $K' =  R(V \to X)_! K$ and hence
$$
\Hom_{D(\mathcal{O}_X)}(K'[n], E) = \Hom_{D(\mathcal{O}_V)}(K[n], E|_V)
$$
Thus by Lemma \ref{lemma-orthogonal-koszul-complex} the vanishing of
these groups implies that $E|_V$ is isomorphic to
$R(U \cap V \to V)_*E|_{U \cap V}$. This implies that $E = R(U \to X)_*E|_U$
(small detail omitted). If this is the case then
$$
\Hom_{D(\mathcal{O}_X)}(Q[n], E) = \Hom_{D(\mathcal{O}_U)}(Q|_U[n], E|_U)
$$
which contains $\Hom_{D(\mathcal{O}_U)}(P[n], E|_U)$ as a direct summand.
Thus by our choice of $P$ the vanishing of these groups implies that $E|_U$
is zero. Whence $E$ is zero.
\end{proof}